\renewcommand{\abstractname}{ABSTRAK}
\begin{abstractind}
Perubahan iklim memberikan dampak signifikan terhadap sektor pertanian, khususnya di Kabupaten Aceh Besar karena merupakan salah satu daerah lumbung pangan Provinsi Aceh. Fluktuasi produksi padi yang tidak menentu memerlukan pengelolaan data iklim yang akurat, konsisten, dan terotomatisasi. Penelitian ini bertujuan untuk merancang dan mengimplementasikan sistem otomasi pembersihan data iklim menggunakan metode \textit{smoothing} data BMKG dan NASA POWER periode 2005\--2025 dengan jumlah 7565 \textit{records} BMKG dan 7578 \textit{records} NASA POWER. Praproses dilakukan melalui kombinasi metode imputasi dan penghalusan data yang disesuaikan dengan karakteristik masing-masing parameter iklim, meliputi \textit{Cubic Spline} untuk suhu dan kelembaban, \textit{Two-Stage Binary Conditional} untuk curah hujan pada data BMKG, serta \textit{Adaptive Exponential Smoothing} untuk mengurangi fluktuasi acak pada data NASA POWER. Hasil praproses menunjukkan bahwa metode yang diterapkan mampu mempertahankan pola dan tren data secara konsisten. Pada data BMKG, parameter suhu dan kelembaban menunjukkan kualitas penghalusan yang sangat baik dengan pelestarian tren di atas 99\%, sedangkan curah hujan mencapai pelestarian tren sebesar 90,12\%. Pada data NASA POWER, parameter suhu menunjukkan kualitas penghalusan yang baik dengan pelestarian tren sebesar 83,5\%, sementara parameter kelembaban tetap mampu mempertahankan tren di atas 82\% meskipun memiliki kualitas penghalusan yang lebih rendah. Hasil penelitian menunjukkan bahwa sistem yang dikembangkan mampu mengotomasi praproses data iklim secara konsisten dan menghasilkan dataset yang siap digunakan pada tahap peramalan untuk mendukung penyusunan kalender tanam digital.

\bigskip
\noindent
\textbf{Kata kunci:} otomasi pemrosesan data, praproses data, smoothing, GCV, pelestarian tren, BMKG, NASA POWER.
\end{abstractind}
